\documentclass[11pt]{article}

\usepackage[margin=1in]{geometry}
\usepackage{amsmath,amssymb,amsthm}
\usepackage{mathtools}
\usepackage{bm}
\usepackage{enumitem}
\usepackage{booktabs}
\usepackage{hyperref}
\usepackage{xcolor}

% ---- notation ----------------------------------------------------------------
\newcommand{\Vocab}{\mathbb{V}}
\newcommand{\Real}{\mathbb{R}}
\newcommand{\Ind}{\mathbb{1}}
\newcommand{\given}{\mid}
\newcommand{\pLM}{p_{\mathrm{LM}}}
\newcommand{\eos}{\textsc{eos}}
\newcommand{\bos}{\textsc{bos}}
\newcommand{\copyact}{\textsf{copy}}
\newcommand{\subact}{\textsf{sub}}
\newcommand{\insact}{\textsf{ins}}
\newcommand{\delact}{\textsf{del}}
\newcommand{\dDL}{d_{\mathrm{DL}}}
\newcommand{\lse}{\operatorname{logsumexp}}
\newcommand{\Cat}{\mathrm{Cat}}
\newcommand{\softmax}{\operatorname{softmax}}
\newcommand{\GFI}{\textsc{gfi}}
\newcommand{\code}[1]{\texttt{#1}}
\newcommand{\chan}{g}            % channel emission score
\newcommand{\wdel}{w_{\delact}}
\newcommand{\wins}{w_{\insact}}

\newtheorem{theorem}{Theorem}
\newtheorem{lemma}{Lemma}
\newtheorem{proposition}{Proposition}
\newtheorem{definition}{Definition}
\newtheorem{remark}{Remark}

% Let TeX stretch interword space a little harder before declaring an overfull line.
\setlength{\emergencystretch}{3em}

\title{A Rao--Blackwellized Pair-HMM Noisy-Channel Model of\\
Language Comprehension and Correction}
\author{genjax noisy-channel port \\ \small \texttt{src/genjax\_port/}}
\date{\today}

\begin{document}
\maketitle

\begin{abstract}
We describe the current noisy-channel model implemented in \code{src/genjax\_port/}. A speaker
draws an \emph{intended} sentence word-by-word from an autoregressive language-model (LM) prior; a
\emph{word-level noisy channel} corrupts it---substituting mistyped words, inserting spurious
words, and deleting intended ones---into the \emph{observed} sentence; comprehension is the
posterior over intended sentences. The central modelling decision is that \emph{the edit alignment
between intended and observed words is never sampled}: it is marginalized analytically by a forward
dynamic program (a pair-HMM run at
the word level, whose per-cell emission is itself a character-level pair-HMM). The only sampled
latent is the intended sentence. Inference is then a Rao--Blackwellized sequential Monte Carlo
(SMC) filter over intended words, with a fully-adapted (channel-aware) proposal whose incremental
weight is a marginal step-likelihood, and a rejuvenation kernel that revisits earlier words as
later context disambiguates them.

The document is in two parts. \textbf{Part~I} is the theoretical framework: the LM prior, the
two-level pair-HMM channel, the Rao--Blackwellized SMC filter, and reanalysis by rejuvenation,
together with the correctness statements (proper weighting, unbiased evidence, posterior
invariance of the moves). \textbf{Part~II} is the engineering that makes it run: vectorization over
particles under GenJAX \code{vmap}/\code{jit}, multi-token intended words, the key/value (KV) cache
that powers rejuvenation, deduplication over the redundant post-resample particle set, and candidate
generation.
Part~I is meant to read independently of Part~II.
\end{abstract}

\tableofcontents

\part{Theoretical framework}
\label{part:theory}

\section{The generative story and the inference target}
\label{sec:overview}

\paragraph{Two levels.} The model is a hierarchy of two sources of structure. At the \emph{sentence}
level, a language model emits a sequence of intended words; at the \emph{character} level, each
intended word may be corrupted into the surface string the reader actually sees. Between them sits
the channel's combinatorial freedom: an intended word may be dropped, and a spurious word may
appear. Writing $x = (x_1,\dots,x_n)$ for the intended sentence (a sequence of \emph{words}, each a
span of one or more subword tokens; see \S\ref{sec:lm}) and $o = (o_1,\dots,o_M)$ for the observed
sentence (segmented
deterministically into $M$ word/punctuation units), the joint is
\begin{equation}\label{eq:joint-abstract}
  p(x, o) \;=\; \underbrace{p_{\mathrm{prior}}(x)}_{\text{LM over intended sentences}}
              \;\cdot\; \underbrace{p_{\mathrm{chan}}(o \given x)}_{\text{word-level noisy channel}} .
\end{equation}
Comprehension/correction is the posterior $p(x \given o) \propto p(x,o)$; the most probable
intended sentence is the model's reconstruction of what the speaker meant.

\paragraph{What is and is not latent.} The intended sentence $x$ is latent. The segmentation of $o$
into units is a deterministic function of $o$ (\code{noise\_word.segment\_words}), so it is
\emph{indexing data}, not a latent variable---every particle consumes the same observed units in
the same order. The \emph{alignment} between intended and observed words---which intended word
produced which observed unit, which observed unit is spurious, which intended word was
dropped---\emph{is} a latent object, but it is one we will never instantiate: Section~\ref{sec:channel}
marginalizes it in closed form. This is the Rao--Blackwellization at the heart of the model.

\section{The language-model prior}
\label{sec:lm}

Let $\Vocab$ be the subword vocabulary of a fixed autoregressive LM with next-token distribution
$\pLM(\cdot \given u)$ for context $u \in \Vocab^{*}$ (the implementation uses Pythia / GPT-NeoX via
penzai; see \S\ref{sec:lm-engine}). The LM works in byte-pair-encoding (BPE) tokens, which split text
into subword pieces, so a single written word may be one token or several. An intended sentence is a
sequence of words $x = (x_1,\dots,x_n)$, where each word $x_j = (t_{j,1},\dots,t_{j,m_j}) \in
\Vocab^{m_j}$ is a span of $m_j \ge 1$ BPE tokens. Concatenating the spans and prefixing a fixed seed $s$ (a
$\bos=\eos$ token, optionally followed by a content-neutral prime; \S\ref{sec:lm-engine}) yields the
token stream the LM scores. The prior factorizes autoregressively over \emph{tokens}, and a complete
sentence is terminated by $\eos$:
\begin{equation}\label{eq:lm-prior}
  p_{\mathrm{prior}}(x)
  \;=\; \Bigg[\prod_{j=1}^{n} \prod_{r=1}^{m_j}
        \pLM\!\big(t_{j,r} \given s, t_{<(j,r)}\big)\Bigg]\;
        \pLM\!\big(\eos \given s, t_{1:n}\big),
\end{equation}
where $t_{<(j,r)}$ is the token prefix preceding the $r$-th token of word $j$. We write the
\emph{word}-level chain-rule factor
\begin{equation}\label{eq:lm-word}
  \ell(x_j \given x_{<j}) \;=\; \sum_{r=1}^{m_j} \log \pLM\!\big(t_{j,r}\given s, x_{<j}, t_{j,1:r-1}\big)
\end{equation}
for the log-probability of emitting word $x_j$ after the words $x_{<j}$. For a single-token word
($m_j = 1$) this is one LM evaluation; multi-token words are the subject of \S\ref{sec:multitoken}.
The sentence length $n$ is itself latent---the model emits words until it draws $\eos$---so the
prior is a distribution over variable-length intended sentences. An intended sentence shorter than
$o$ is the prior's way of positing spurious (inserted) observed words; one longer than $o$ posits
dropped (deleted) intended words. \emph{Tempering.} It is sometimes useful to flatten the LM's
(over-confident) preferences. We therefore allow a temperature $\lambda \in (0,1]$ on the prior,
targeting $p_{\mathrm{prior}}^{\lambda}\,p_{\mathrm{chan}}$; $\lambda = 1$ is the untempered model
(\S\ref{sec:smc} treats $\lambda$ as a target choice).

\section{The noisy channel as a two-level pair-HMM}
\label{sec:channel}

The channel likelihood $p_{\mathrm{chan}}(o \given x)$ must account for three edit types between the
intended word sequence $x_{1:n}$ and the observed unit sequence $o_{1:M}$:
\begin{itemize}[leftmargin=1.6em,itemsep=1pt]
  \item \textbf{copy / substitute} --- an intended word $x_j$ is realized as an observed unit
        $o_k$, exactly (copy) or with character-level corruption (substitution);
  \item \textbf{insertion} --- an observed unit $o_k$ is \emph{spurious}: it corresponds to no
        intended word;
  \item \textbf{deletion} --- an intended word $x_j$ is \emph{missing}: it produces no observed
        unit.
\end{itemize}
A particular assignment of these operations is an \emph{alignment} $\mathcal A$. The channel
likelihood marginalizes over all alignments,
\begin{equation}\label{eq:chan-marginal}
  p_{\mathrm{chan}}(o \given x) \;=\; \sum_{\mathcal A} p(o, \mathcal A \given x),
\end{equation}
and the contribution of an alignment is a product of local factors: a character-level emission score
for each copy/substitute pairing, a cost $\wins$ for each spurious unit, and a cost $\wdel$ for each
missing word. The key structural fact is that \eqref{eq:chan-marginal} is computed by a forward
dynamic program---a pair-HMM at the word level---so it never enumerates alignments explicitly.

\subsection{Inner level: the character channel}
\label{sec:char-channel}

The emission score of explaining an observed unit $o_k$ by an intended word $x$ is the
log-likelihood of a character-level edit channel between their \emph{surface forms}---the written
character strings $\mathrm{surf}(\cdot)$ of the two words:
\begin{equation}\label{eq:char-channel}
  \chan(o_k \given x) \;=\; \log P_{\mathrm{char}}\!\big(\mathrm{surf}(o_k) \given \mathrm{surf}(x)\big).
\end{equation}
This is itself a forward pair-HMM over characters with copy, substitute, insert, delete, and
adjacent-transposition edges (\code{channel\_logpdf}). Writing
$G[i][j] = \log P(\mathrm{obs}[:j]\given \text{first } i \text{ intended chars})$, the recurrence
combines four arcs at each cell---copy-or-substitute a character, delete a character (the intended
character produces nothing), insert a character (an observed character with no intended source), and
transpose a pair of adjacent characters:
\begin{equation}\label{eq:char-dp}
\begin{split}
  G[i][j] = \lse\big(\, & G[i{-}1][j{-}1] + c_{ij},\ \ G[i{-}1][j] + \log\theta_{\delact}, \\
                        & G[i][j{-}1] + \log\theta_{\insact},\ \ G[i{-}2][j{-}2] + \log\theta_{\mathrm{tr}}\,\big),
\end{split}
\end{equation}
where $c_{ij} = \log\theta_{\copyact}$ if the two characters match and $\log\theta_{\subact}$
otherwise, and the transposition arc fires only when the trailing two characters are swapped. The
default character rates are $\theta_{\copyact} = 0.9$, $\theta_{\subact} = (1-0.9)/26$,
$\theta_{\insact}=\theta_{\delact}=0.05$, $\theta_{\mathrm{tr}}=\theta_{\subact}$. Reading
$G[n_x][n_o]$ (all characters consumed on both sides) gives \eqref{eq:char-channel}. A \emph{copy}
($x$ has the same surface form as $o_k$) scores near $0$; a substitution whose surface forms differ by
$d$ character edits scores roughly $d\cdot\log\theta_{\subact}$. Here ``$d$ character edits'' is the
Damerau--Levenshtein distance---the smallest number of single-character insertions, deletions,
substitutions, and adjacent transpositions that turns one string into the other---the same metric the
candidate generator uses (\S\ref{sec:candidates}). We write $\chan(o_k\given x)$ for the resulting per-pairing
emission score and assemble, for a sentence, the emission table
$\mathrm{EMIT}[k, c] = \chan(o_k \given c)$ over the candidate intended words $c$.

\subsection{Outer level: the word channel as a forward DP}
\label{sec:word-channel}

Lift the same three-way recurrence one level, from characters to words. A particle that has emitted
intended words $x_{1:n}$ carries a \emph{forward vector} $\alpha^{(n)} \in (\Real\cup\{-\infty\})^{M+1}$,
\begin{equation}\label{eq:alpha-def}
  \alpha^{(n)}[k] \;=\; \log P\big(o_{1:k},\ \text{exactly } k \text{ observed units consumed}
                                   \;\big|\; x_{1:n}\big),
\end{equation}
i.e.\ the total channel mass (summed over alignments) for the hypothesis that the first $n$ intended
words have produced exactly the first $k$ observed units. This vector \emph{is} the
Rao--Blackwellized alignment: it is a fixed-length ($M{+}1$) sufficient statistic, deterministic
given $x_{1:n}$ and $o$, that replaces an explicit sampled alignment.

Emitting one more intended word $x_{n+1}$ updates the vector by one pair-HMM row
(\code{\_word\_row\_update}). With emission column
$\mathrm{e}[k] = \chan(o_{k}\given x_{n+1})$ and using $\beta$ as the intermediate after the
copy/substitute and missing-word arcs,
\begin{align}
  \beta[0]   &= \alpha^{(n)}[0] + \wdel, \label{eq:row-beta0}\\
  \beta[k]   &= \lse\big(\underbrace{\alpha^{(n)}[k{-}1] + \mathrm{e}[k]}_{\text{align } x_{n+1}\to o_k},\;
                          \underbrace{\alpha^{(n)}[k] + \wdel}_{x_{n+1}\text{ missing}}\big),
                          \quad k\ge 1, \label{eq:row-beta}\\
  \alpha^{(n+1)}[k] &= \lse\big(\beta[k],\;\alpha^{(n+1)}[k{-}1] + \wins[k]\big),
                          \quad (\text{$o_k$ spurious}) . \label{eq:row-ins}
\end{align}
Equation~\eqref{eq:row-beta} combines the diagonal arc (align $x_{n+1}$ to $o_k$: substitute or copy)
and the up arc (the intended word $x_{n+1}$ is missing, cost $\wdel$). Equation~\eqref{eq:row-ins} is
a left-to-right pass that lets a spurious observed unit $o_k$ advance the consumption count at cost
$\wins[k]$ without consuming any intended word. The vector is initialized at $\alpha^{(0)}[k] = \sum_{j\le k}
\wins[j]$ (the first $k$ observed units are all spurious, i.e.\ leading insertions), and the
\emph{channel marginal} of a completed intended sentence is the full-consumption entry
\begin{equation}\label{eq:chan-from-alpha}
  p_{\mathrm{chan}}(o \given x_{1:n}) \;=\; \exp\,\alpha^{(n)}[M].
\end{equation}
The partial forward mass $\lse_k \alpha^{(n)}[k]$ is the likelihood of the \emph{observed-prefix-so-far}
under the partial intended sentence---the quantity the SMC filter tracks before the sentence is
complete (\S\ref{sec:smc}).

\paragraph{Edit costs.} $\wdel$ (missing-word, the \emph{over-editing} knob; default $-9$ nats) is
the price of adding a word to the reconstruction that the observation does not support; too cheap and
the filter hallucinates fluent words, too steep and it stops restoring genuinely dropped words.
$\wins[k]$ (spurious-word) is by default \emph{frequency-aware}: $\wins[k] = \log\rho_{\insact} -
\mathrm{surp}_{\mathrm{uni}}(o_k)$, a per-unit decomposition into an insertion rate $\rho_{\insact}$
(default $0.02$) times the out-of-context unigram content distribution. Making the cost depend on the
word's base frequency keeps a rare-but-correct word from being dropped as spurious: under a flat cost
(say $-\log V$, uniform over the vocabulary), any word rarer than uniform would be cheaper to discard
than to keep, whereas here a rare word is expensive to explain away (\code{unigram.py}; the cost
enters \eqref{eq:row-ins} as a per-observed-word vector).

\paragraph{The band: bounded reach for insertions and deletions.} Insertions and deletions let the
intended-word count $n$ and
the consumed-unit count $k$ drift apart. To keep the DP from positing arbitrarily many consecutive
spurious or missing words, the forward vector is masked to a diagonal band
$|k - t| \le b$ after the $t$-th word step (default $b = 2$; \code{band\_mask}). The band is what
gives insertions and deletions their bounded reach; it is the only approximation in the otherwise
exact channel marginal, and at $b \ge M$ it is vacuous.

\begin{remark}[No explicit edit action]
There is deliberately \emph{no} sampled ``edit type.'' Substitution and deletion are the diagonal and
up arcs of \eqref{eq:row-beta}; insertion is the left-to-right pass of \eqref{eq:row-ins}. Scoring a
spurious word as if it were a word the LM chose to emit would be a category error---a spurious word is
a channel event, not an LM emission---and would double-count its probability mass, biasing the
evidence by $\approx 0.2$ nats against exact enumeration. Every SMC step is thus a clean choice of the
next intended word (or $\eos$); the edits live entirely inside the marginalized channel.
\end{remark}

\section{Rao--Blackwellized sequential Monte Carlo}
\label{sec:smc}

\subsection{Sequence of targets}

Combining \eqref{eq:lm-prior} and \eqref{eq:chan-from-alpha}, the (tempered) unnormalized posterior
over \emph{complete} intended sentences is
\begin{equation}\label{eq:final-target}
  \gamma(x_{1:n}) \;=\; \Big[\textstyle\prod_{j} \pLM^{\lambda}(x_j\given x_{<j})\Big]\,
                        \pLM^{\lambda}(\eos\given x_{1:n})\;\exp\alpha^{(n)}[M],
\end{equation}
with normalizer $Z = \sum_{x} \gamma(x) = p(o)$ (up to the deterministic constant
$\lse_k\alpha^{(0)}[k]$ folded out by the implementation). The filter realizes this as a sequence of
intermediate targets indexed by the number of emitted intended words. After $n$ words the
\emph{partial} unnormalized target is
\begin{equation}\label{eq:interm-target}
  \tilde\gamma_n(x_{1:n}) \;=\; \Big[\textstyle\prod_{j\le n} \pLM^{\lambda}(x_j\given x_{<j})\Big]\;
                        \exp\Big(\lse_k \alpha^{(n)}[k]\Big),
\end{equation}
the prior of the first $n$ words times the likelihood of the observed-prefix-so-far. A particle is
extended word-by-word under \eqref{eq:interm-target}; when it draws $\eos$ it commits, and a
\emph{terminal correction} converts the partial mass into the full-consumption marginal
\eqref{eq:chan-from-alpha}:
\begin{equation}\label{eq:term-correction}
  \Delta^{\mathrm{term}} \;=\; \alpha^{(n)}[M] - \lse_k \alpha^{(n)}[k].
\end{equation}

\subsection{The channel-aware (fully-adapted) proposal}
\label{sec:proposal}

At each step a live particle proposes its next intended word from a finite \emph{candidate set}
$\mathcal C$ assembled from two sources (\code{\_caprop\_scores}):
\begin{enumerate}[leftmargin=1.8em,itemsep=1pt]
  \item \textbf{channel-compatible words}: the emission candidates of the observed units in a small
        window around the current alignment frontier $\arg\max_k\alpha^{(n)}[k]$---the words that
        could plausibly have produced the next observed unit (a copy of it, or a near substitution);
  \item \textbf{fluency words}: the top-$J$ words under the live LM $\pLM(\cdot\given x_{\le n})$.
        These supply candidates for deletions (a dropped word the LM expects to see) and ensure the
        proposal always has plausible options even when the observation is uninformative.
\end{enumerate}
Candidates that share a surface form are merged (kept once), and $\eos$ is always a candidate. Each
candidate word $c$ is scored by the increment it would add to the intermediate log-target
\eqref{eq:interm-target}:
\begin{align}
  S(c)    &\;=\; \lambda\,\ell(c\given x_{\le n})
            \;+\; \underbrace{\Big[\lse_k \alpha'_c[k] - \lse_k\alpha^{(n)}[k]\Big]}_{
                  \Delta_n(c):\ \text{gain in channel forward mass}}, \label{eq:cand-score}\\
  S(\eos) &\;=\; \lambda\log\pLM(\eos\given x_{\le n}) + \Delta^{\mathrm{term}}, \label{eq:cand-score-eos}
\end{align}
where $\alpha'_c = \texttt{band\_mask}(\texttt{row\_update}(\alpha^{(n)}, \mathrm{EMIT}[:,c]))$ is the
forward vector after emitting $c$. The proposal is the softmax of these scores,
\begin{equation}\label{eq:proposal}
  q_n(c) \;=\; \softmax_{c\in\mathcal C\cup\{\eos\}} S(c),
\end{equation}
and the next word is drawn $x_{n+1} \sim q_n$. This is the locally optimal (``fully-adapted'')
proposal: it scores each candidate by exactly its contribution to the next target, so the
information from \emph{both} the LM and the channel shapes the draw.

\begin{lemma}[Incremental weight collapse]\label{lem:collapse}
Under the proposal \eqref{eq:proposal}, the incremental SMC importance weight is independent of the
sampled word and equals the log-sum of candidate scores:
\begin{equation}\label{eq:wstep}
  \log w_{n+1} \;=\; \frac{\tilde\gamma_{n+1}(x_{1:n+1})}{\tilde\gamma_n(x_{1:n})\,q_n(x_{n+1})}
               \;=\; \lse_{c\in\mathcal C\cup\{\eos\}} S(c).
\end{equation}
\end{lemma}
\begin{proof}
The target increment for choosing word $c$ is exactly $\exp S(c)$ by construction
(\eqref{eq:cand-score} is the log-ratio $\tilde\gamma_{n+1}/\tilde\gamma_n$ at $x_{n+1}=c$). With
$q_n(c) = \exp S(c)/\sum_{c'}\exp S(c')$, the ratio in \eqref{eq:wstep} is
$\exp S(c)\big/\!\big(\exp S(c)/\sum_{c'}\exp S(c')\big) = \sum_{c'}\exp S(c')$. The chosen-word
factor cancels; taking logs gives the $\lse$. \qed
\end{proof}
This is the line \code{incr = logsumexp(scores)} in \code{\_make\_kernel}; the weight is the marginal
step-likelihood under the candidate set, and it is folded into the running log-evidence.

\begin{remark}[Candidate-restricted support]
The candidate set $\mathcal C$ is finite and data-driven, so the filter targets the posterior
\emph{restricted to reachable intended sentences} (those expressible from window + top-$J$ + COPY
candidates). On the toy bigram model the candidate set is the entire vocabulary, so the restriction
is vacuous and the filter is exact---this is what the brute-force certification in
\code{tests/test\_pairhmm\_exact.py} checks (\S\ref{sec:certification}). On Pythia the set is a
high-recall pool; words outside it are unreachable, which is the intended behaviour (a correction
must be a near neighbour or an LM-plausible bridge).
\end{remark}

\subsection{Resampling and the algorithm}

Particle weights are accumulated in log space and resampling is triggered adaptively when the
effective sample size---a standard measure of how many particles carry non-negligible
weight---falls below $P/2$ (multinomial resampling via \code{jax.random.categorical} on the
normalized log-weights). Algorithm~\ref{alg:smc} is one run; ``slack'' is a small number of extra word
steps beyond $M$, which give inferred deletions room to make the intended sentence longer than the
observed one.

\begin{figure}[t]
\hrule\vspace{3pt}
\textbf{Algorithm 1: Rao--Blackwellized pair-HMM SMC.}
\begin{enumerate}[leftmargin=2em,itemsep=1pt]
  \item Seed each of $P$ particles: context buffer $\gets s$, $\alpha \gets \alpha^{(0)}$,
        $\mathrm{done}\gets\textsf{false}$, $\log w \gets 0$, $\log\widehat Z\gets 0$.
  \item For step $n = 0,1,\dots,M+\text{slack}-1$:
  \begin{enumerate}[leftmargin=1.4em,itemsep=0pt]
    \item One batched LM forward over all particles $\to \pLM(\cdot\given x_{\le n})$.
    \item Assemble candidate scores $S(\cdot)$ \eqref{eq:cand-score}; sample $x_{n+1}\sim q_n$
          \eqref{eq:proposal} (live particles only; $\mathrm{done}$ particles idle).
    \item $\log w \mathrel{+}= \lse_c S(c)$ (Lemma~\ref{lem:collapse}); append the word's token span
          to the buffer; update $\alpha$ \eqref{eq:row-beta}--\eqref{eq:row-ins}; set
          $\mathrm{done}$ on $\eos$.
    \item If $\mathrm{ESS}(\log w) < P/2$: fold $\log\widehat Z\mathrel{+}= \lse_p\log w^{(p)} -\log P$;
          resample ancestors; reset $\log w\gets 0$;
          \emph{(optionally) run a rejuvenation sweep} (\S\ref{sec:rejuv}) and fold its weight.
  \end{enumerate}
  \item Apply the terminal correction \eqref{eq:term-correction} to live particles; fold the final
        $\lse_p\log w^{(p)}-\log P$ into $\log\widehat Z$.
  \item Return the weighted particles (posterior over intended sentences) and $\log\widehat Z$.
\end{enumerate}
\hrule
\label{alg:smc}
\end{figure}

\subsection{Correctness}

\begin{definition}[Proper weighting]
A weighted set $\{(x^{(p)}, w^{(p)})\}_{p=1}^P$ is \emph{properly weighted} for an unnormalized
target $\gamma$ with $Z=\sum\gamma$ if $\mathbb E[\tfrac1P\sum_p w^{(p)} f(x^{(p)})] = \sum f\,\gamma
= Z\,\mathbb E_{\bar\gamma}[f]$ for all integrable $f$.
\end{definition}

\begin{theorem}[Properly weighted filter; unbiased evidence]\label{thm:smc}
For every $n$, the weighted particles after step $n$ of Algorithm~\ref{alg:smc} are properly weighted
for $\tilde\gamma_n$ \eqref{eq:interm-target}; after the terminal correction they are properly
weighted for $\gamma$ \eqref{eq:final-target}. Consequently $\widehat Z = \prod (\tfrac1P\sum_p
w^{(p)})$ is unbiased for the marginal likelihood, $\mathbb E[\widehat Z] = Z = p(o)$.
\end{theorem}
\begin{proof}[Proof sketch]
Induction on $n$, identical in structure to standard sequential importance sampling with resampling.
The base case is the seeded particle set (weight $1$, $\alpha^{(0)}$). The inductive step is one
importance-sampling extension with proposal $q_n$ \eqref{eq:proposal} and target increment
$\exp S(\cdot)$: because $q_n>0$ on the
support of the increment (the proposal is the softmax of the same scores), the importance-sampling
identity gives proper weighting for $\tilde\gamma_{n+1}$, with incremental weight \eqref{eq:wstep}. Multinomial
resampling replaces weights by their mean and is unbiased, and folding the mean into $\widehat Z$
keeps the product unbiased by the tower rule. The terminal correction multiplies each live particle's
weight by $\exp\Delta^{\mathrm{term}}$, converting $\tilde\gamma_n \to \gamma$ pointwise. The
Rao--Blackwellization is exact: $\alpha^{(n)}$ in \eqref{eq:alpha-def} equals the channel mass
summed over all alignments \eqref{eq:chan-marginal}, so replacing a sampled alignment by $\alpha$
loses no target mass and (strictly) reduces weight variance. \qed
\end{proof}

\begin{remark}[Why Rao--Blackwellize]
Sampling the alignment would make each particle a guess at both the intended sentence and its
edit structure; the weights would carry the variance of that second guess. Summing the alignment out
means a particle is judged only on its intended sentence, with the channel's verdict computed
exactly. The fully-adapted proposal then makes the per-step weight a near-deterministic marginal
likelihood (Lemma~\ref{lem:collapse}), so small particle counts suffice ($P\!\approx\!64$--$256$ on
Pythia).
\end{remark}

\section{Reanalysis by rejuvenation}
\label{sec:rejuv}

\subsection{Why}
Resampling repeatedly prunes alternatives. A word committed under local information early in the
filtering pass can never be revised once \emph{later} context disambiguates it---the classic
garden-path situation. Rejuvenation interleaves, after a resampling step, Markov-chain Monte Carlo
(MCMC) moves that revisit recent words within a lookback window. These moves target the \emph{same}
posterior; they do not change the target, they let the chain reach high-posterior reanalyses and
restore diversity to a particle set that resampling has collapsed onto a few distinct hypotheses
(\code{pairhmm\_rejuv.py}).

\subsection{The full-conditional (Gibbs / SMCP3) word move}
\label{sec:rejuv-move}

A sweep resamples one intended word at a time, in place, from its full conditional. For word slot
$w$ in the window, holding the other words fixed, the conditional is
\begin{equation}\label{eq:full-cond}
  P(x_w = c \given x_{\setminus w}, o)
  \;\propto\; \underbrace{p_{\mathrm{prior}}^{\lambda}(x_{1:n}\text{ with } x_w{=}c)}_{\text{LM of the spliced sentence}}
              \;\cdot\; \underbrace{\exp\alpha[M]\ \text{(or }\lse_k\alpha[k])}_{\text{channel marginal, recomputed}},
\end{equation}
over a candidate set $\mathcal C_w = \{\copyact\}\cup \mathcal P_w$ (keep the current word, or a
SymSpell pool entry; \S\ref{sec:candidates}). Both factors are scored by exactly the pieces the
forward filter uses---the autoregressive LM for the prior, the word-level DP for the channel
marginal---so sampling from $\softmax_{c} \log[\,\cdot\,]$ \emph{is} a draw from the conditional.
``Done'' (complete-hypothesis) particles score the terminal channel marginal $\alpha[M]$ plus an
$\eos$ term; mid-loop particles score the partial forward mass $\lse_k\alpha[k]$.

The move is consumed in either of two equivalent ways, both built on GenJAX's \code{Rejuvenate}
SMCP3 recipe (\code{inference/requests/rejuvenate.py}; the move proposes from the conditional,
\code{Update}s the trace, and assesses the reverse, returning the SMCP3 weight
$\mathcal W = w_{\mathrm{upd}} + s_{\mathrm{bwd}} - s_{\mathrm{fwd}}$; the algebra is recalled in
Appendix~\ref{app:gfi}):
\begin{description}[leftmargin=1.4em]
  \item[(A) Metropolis--Hastings.] Accept the move with probability $\min(1, e^{\mathcal W})$;
        weights untouched.
  \item[(B) SMCP3 reweighting.] Move deterministically and multiply the particle weight by
        $e^{\mathcal W}$. The implementation uses (B): it applies the sweep to the equal-weight
        post-resample cloud and folds $\mathcal W$ into $\log w$ \emph{before} the next resample, so
        an asymmetric candidate set carries real mass into the resample.
\end{description}
For a true full-conditional proposal the SMCP3 weight is $\mathcal W \approx 0$ (asserted as a
built-it-right check); the weight becomes informative exactly when the candidate set is asymmetric
(e.g.\ COPY plus a one-sided neighbour pool).

\begin{theorem}[Invariance]\label{thm:rejuv}
Each word move leaves the posterior $\bar\gamma\propto p(x\given o)$ invariant: consumed as (A) it
satisfies detailed balance w.r.t.\ $\bar\gamma$; consumed as (B) it preserves proper weighting for
$\bar\gamma$. A windowed sweep---any data-dependent but state-independent schedule of such moves---is
therefore also $\bar\gamma$-invariant.
\end{theorem}
\begin{proof}[Proof sketch]
The move overwrites slot $w$ and leaves all other addresses fixed, so the forward proposal of the
new word and the reverse proposal of the old word are both scored by the same conditional kernel
evaluated at the respective points. (A) is then the standard Metropolis--Hastings acceptance for
proposal $Q$ and target $\bar\gamma$, with $\mathcal W$ the log of the MH ratio (the \code{Update}
weight supplies the target ratio, the backward request supplies the addresses the reverse kernel
scores); detailed balance is the textbook identity. (B) is the SMCP3 theorem specialized to a
deterministic overwrite with auxiliary randomness from $Q$: replacing the proper-weighting integral
and changing variables $(x,\zeta)\mapsto x'$ cancels the $Q$ factors and preserves
$\sum f\,\bar\gamma$. The schedule invariance is immediate: a composition of $\bar\gamma$-invariant
kernels is $\bar\gamma$-invariant regardless of how the (data-dependent) schedule chooses them. \qed
\end{proof}

\subsection{Scheduling: the surprisal-relative gate}
\label{sec:gate}

\emph{Where} to spend rejuvenation is decided by a gate on each word's surprisal (its negative
log-probability). Gating on raw contextual surprisal fires on any word that is improbable in
context---including words that are merely \emph{rare} (a rare word is \emph{expected} to be
surprising, so revisiting it wastes effort and can pull a legitimately rare literal toward a common
substitute). Instead the gate fires on contextual surprisal \emph{relative to} the word's
out-of-context (unigram) surprisal:
\begin{equation}\label{eq:gate}
  p_{\mathrm{rejuv}}(o_k) \;=\; \sigma\!\big(\beta\,[\,\mathrm{surp}_{\mathrm{ctx}}(o_k)
                                              - \mathrm{surp}_{\mathrm{uni}}(o_k) - \tau\,]\big),
\end{equation}
which fires only when a word is \emph{more} surprising in context than its base rate predicts---the
signature of a genuine garden-path reanalysis rather than a rare word (\code{unigram.py}). Because
the gate is data-dependent but state-independent, Theorem~\ref{thm:rejuv} guarantees it changes only
mixing, not the stationary distribution.

\subsection{The cost of reanalysis}
The LM is causal: overwriting an intended word at slot $w$ changes $\pLM$ for all later positions.
Scoring the conditional \eqref{eq:full-cond} therefore re-scores the suffix from $w$ to the frontier;
the \emph{prefix cancels} in the softmax (it is identical across candidates). A lookback window of
$\ell$ words bounds the suffix at $O(\ell)$ LM evaluations per candidate per move---and re-scoring
the suffix is exactly what lets later context vote on the earlier word. Computing that suffix
efficiently is the central engineering problem of Part~II (\S\ref{sec:kv}).

\part{Engineering and implementation}
\label{part:eng}

Part~I is a specification; Part~II is what makes it run at usable speed. The throughline is that the
LM forward pass dominates cost, so every optimization is about \emph{reducing or sharing} LM work:
batch it across particles (\S\ref{sec:vmap}), avoid recomputing the prefix in rejuvenation
(\S\ref{sec:kv}), and never score a buffer twice (\S\ref{sec:dedup}). Multi-token words
(\S\ref{sec:multitoken}) and candidate generation (\S\ref{sec:candidates}) are the two places where
the BPE tokenizer's reality intrudes on the clean word-level story of Part~I.

\section{Vectorization over particles (GenJAX \texorpdfstring{\code{vmap}}{vmap}/\texorpdfstring{\code{jit}}{jit})}
\label{sec:vmap}

\paragraph{One filter, two models.} All inference lives in \code{pairhmm\_smc.py} and is generic; the
model-specific pieces---the LM, the channel, the candidate generator, the seed---are injected through
a \code{PairHMMModel} dataclass. The toy bigram (fast, no LM load, certified by brute force) and
Pythia (\code{pythia\_word\_caprop.py}) are two configurations of the \emph{same} code, so correctness
proven on the toy transfers to Pythia by construction (\S\ref{sec:certification}).

\paragraph{Fixed-shape per-particle state.} Each particle carries a fixed-shape tuple of fields: a
padded token buffer \code{ctx\_buf} and its fill length \code{ctx\_len}; the intended-word count
\code{n\_words}; the per-word token counts \code{word\_len} and channel surface ids \code{word\_surf};
the forward vector $\alpha$; and a completion flag \code{done}. Nothing is ragged, so the whole set of
$P$ particles---the \emph{particle cloud}---is a collection of arrays with a leading axis of size $P$,
and one SMC step is a single \code{jax.vmap} over that axis. Crucially, the \emph{whole cloud shares
one batched LM forward per step}: a single call evaluates $\pLM(\cdot\given x_{\le n})$ for all $P$
particles at once, so the dominant cost is paid once for the whole cloud, which is why runtime is
nearly flat in $P$.

\paragraph{The proposal kernel is a GenJAX \code{@gen} function.} The per-step kernel
(\code{\_make\_kernel}) samples \code{genjax.categorical(scores) @ "action"} and injects the
incremental weight through a deterministic \code{factor(incr) @ "ev"}; the filter drives it with
\code{kernel.importance}, so the SMC weight is the generative-function-interface importance weight
(Appendix~\ref{app:gfi}), not a hand-rolled quantity. The hand-written outer loop and the generative
model thus agree by construction.

\paragraph{Compile once, reuse across runs.} The jitted \code{extend} step and the rejuvenation step
are built by factories memoized on their \emph{structural} signature (sequence lengths, candidate
counts), with all per-run-varying data (the emission table, $\alpha^{(0)}$, the candidate pools, the
edit costs, the seed) threaded as \emph{traced arguments} rather than baked into the compiled
function. Two runs of the same sentence \emph{shape} therefore hit JAX's compilation cache and reuse
the existing compiled code; only a new sentence length triggers a fresh compile. This matters most on
the rejuvenation path, where the step is reused across the many resampling events of a single run as
well as across separate runs.

\section{Multi-token intended words}
\label{sec:multitoken}

The clean ``one word'' of Part~I is, to the BPE tokenizer, a span of one or more tokens. Common words
are usually a single token, but corrections sometimes are not (a misspelling of a rarer word, a
proper noun, morphology), so a candidate intended word is represented as a \textbf{(token span,
surface id)} pair, not a single token.

\paragraph{Decoupling LM scoring from the channel column.} The token span drives the LM chain-rule
\eqref{eq:lm-word} and the buffer splice; the surface id indexes the channel emission table. A
single-token candidate keeps $\text{surface id} = \text{token id}$ and the cheap one-forward LM score;
a multi-token candidate gets its LM score from the chain-rule scorer (\code{tail\_logprobs}, the same
KV-cached machinery as \S\ref{sec:kv}) and its channel column from a per-sentence \emph{augmented}
emission table (multi-token surfaces occupy appended columns). The candidate inventory builder
(\code{\_build\_candidates}) returns the per-word candidate spans, surfaces, and a multi-token index;
$T_{\max}$ is the maximum span length over the sentence.

\paragraph{Capacity-parametric state and the re-pack.} The intended sentence is stored as an array
\code{word\_tok} of shape $[P, W_{\max}, T_{\max}]$ together with per-word lengths, and the flat LM
buffer is rebuilt from it by a cumulative-sum/scatter ``pack'' (\code{\_pack}) that handles unequal
spans collision-free and at fixed shape; \code{\_unpack} inverts it using the stored per-word token
counts. The band uses the \emph{word} count \code{n\_words}, not the token count \code{ctx\_len},
since the two diverge once words are multi-token.

\paragraph{Bit-identical reduction.} The decisive invariant: with no multi-token candidates
($n_{\mathrm{mt}} = 0$, $T_{\max}=1$) every code path and value is \emph{bit-identical} to the
certified single-token filter---each word's single token sits at the first slot of \code{word\_tok},
\code{\_pack} with all lengths $1$ is the identity map from word index to buffer position, and the
span is just the one token. The exact-enumeration gates (\S\ref{sec:certification}) run in this regime and so
continue to certify the multi-token code.

\section{The KV-cache for rejuvenation}
\label{sec:kv}

\subsection{The LM engine}
\label{sec:lm-engine}
The LM is a Pythia (GPT-NeoX) model loaded from HuggingFace and converted \emph{once} to penzai
(\code{lm\_penzai.py}); after conversion the forward pass is pure penzai/JAX (no torch in the loop),
so it \code{jit}s and \code{vmap}s. Position $0$ of every buffer is seeded with $\eos$ as a
start-of-sequence context; an optional content-neutral prime (a ``\code{.}'' that has no trailing
space, to keep inter-word spacing consistent with the leading-space word tokens) conditions the
model mid-document rather than at a document boundary. The LM is exposed only behind
\code{next\_token\_logprobs}, so a different model could be swapped in without touching the filter.

\subsection{Prefix cancellation and the suffix-tail scorer}
The rejuvenation conditional \eqref{eq:full-cond} re-scores a candidate-dependent suffix, but the
prefix LM (the words before the edited slot) is identical for every candidate and \emph{cancels} in
the softmax and the SMCP3 weight. The scorer therefore computes only $\log P(\text{tail}\given
\text{prefix})$ for the tail $=$ [candidate span, suffix words, $\eos$?], over the $K$ candidates of a
slot. A transformer caches, for each layer, the key and value tensors it has already computed for the
context (its ``KV cache''), so that extending the context reuses that work. The scorer
(\code{\_batch\_tail\_logprobs\_kv}) exploits this: it \emph{prefills} the prefix once---runs the
model over the prefix a single time to populate the cache---and then shares that cache across all $K$
candidate feeds. It reads $\log P(\text{tail}_0\given \text{ctx})$ from the prefill output at
position \code{ctx\_len}$-1$, then feeds each short candidate tail starting at position
\code{ctx\_len} and reads the remaining chain-rule terms. The penzai
\code{KVCachingTransformerLM} is driven by an unbind/bind pattern---\code{unbind\_variables} freezes
the shared prefill K/V, and each candidate \code{bind}s a fresh copy and advances
\code{cache\_end\_index}---so the expensive prefix attention is computed exactly once per particle per
word. This is what brings the rejuvenation LM cost down from the naive ``re-score the whole sentence
per candidate'' (an $O(\text{length})$ forward for each of $P\cdot K$ buffers) to one shared prefill
plus cheap tail feeds.

\section{Deduplication over the degenerate cloud}
\label{sec:dedup}

Resample-every-word makes the post-resample cloud highly degenerate---empirically $\sim\!93\%$ of the
$P$ particles share an identical intended prefix (\code{cache\_dedup.py}). Two LM workloads exploit
this, both \emph{exact} (bit-identical posterior given the same RNG; they remove only redundant
forwards):
\begin{itemize}[leftmargin=1.6em,itemsep=1pt]
  \item \textbf{Filter forward (\code{make\_forward\_dedup}).} The per-step batched LM forward is run
        on the \emph{unique} buffers only, then scattered back to all $P$ particles.
  \item \textbf{Rejuvenation tail (\code{\_dedup\_tail}).} The dominant single-sentence cost is the
        sweep's KV prefills, which scale with $P$. The tail scorer is keyed on exactly the bytes it
        reads---the prefix it prefills plus the candidate tails---and run on the unique rows only;
        the per-particle SMCP3 \emph{sample} still runs on all $P$ (duplicates must diverge---that is
        the diversification the dedup must never collapse).
\end{itemize}
Unique rows are padded to a fixed ``bucket ladder'' of batch sizes so the jitted scorer recompiles at
only a few shapes.

\section{Candidate generation}
\label{sec:candidates}

Candidates are retrieved by a \textbf{SymSpell deletion index} over the word-initial vocabulary
(\code{noise\_word.py}). Two strings are within Damerau--Levenshtein distance $D$ iff they share a
common delete-variant (up to $D$ deletions on each side), so indexing the dictionary's delete-variants
lets a query retrieve all candidates within $D$ by looking up its own deletes---the edit distance is a
\emph{real} parameter, not capped at $1$. The verified distance is then computed by an
optimal-string-alignment Damerau--Levenshtein with early cutoff, and candidates are returned
nearest-first, truncated to \code{MAX\_SUB\_CANDIDATES}. The truncation is a tractability bound, not a
modelling limit: far candidates are already crushed by $\theta_{\subact}^{d}$ in the channel
\eqref{eq:char-channel}. There are two pools: the $\sim\!27$k single-token word-initial vocabulary,
and a multi-token pool built from the top-$N$ \code{wordfreq} English words that tokenize to $\ge 2$
BPE tokens (\S\ref{sec:multitoken}). The COPY candidate---the observed unit's own token span,
verbatim---is always first, so a correctly-spelled (or correctly-punctuated) word can always be
emitted as itself; using the \emph{observed} span keeps punctuation and casing faithful (a terminal
``\code{.}'' stays the attached token, not a spaced one).

\section{Reference parameters}
\label{sec:params}

\begin{center}
\small
\begin{tabular}{@{}llll@{}}
\toprule
Symbol & Meaning & Default & Where \\
\midrule
$\theta_{\copyact}$ & char copy probability & $0.90$ & \code{pythia\_word\_caprop} \\
$\theta_{\subact}$  & char substitution prob.\ & $(1{-}0.9)/26$ & \code{pythia\_word\_caprop} \\
$\theta_{\insact},\theta_{\delact}$ & char insert/delete prob.\ & $0.05$ & \code{pythia\_word\_caprop} \\
$\wdel$ & missing-word (over-editing) cost & $-9$ nats & \code{WDEL\_DEFAULT} \\
$\rho_{\insact}$ & spurious-word insertion rate & $0.02$ & \code{ins\_rate} \\
$b$ & alignment band half-width & $2$ & \code{band} \\
$\lambda$ & LM-prior temperature & $1.0$ & \code{lm\_temp} \\
$J$ & top-$J$ LM bridge candidates & $8$ & \code{J} \\
$K_e$ & candidate words per observed word & $12$ & \code{Ke} \\
$\ell$ & rejuvenation lookback window & $3$ words & \code{rejuv\_lookback} \\
$P$ & particle count (Pythia) & $64$--$256$ & \code{P} \\
\bottomrule
\end{tabular}
\end{center}

\noindent The three over-/under-editing knobs ($\wdel$, $\rho_{\insact}$, $\lambda$) interact and are
co-tuned on a gold set: $\lambda<1$ trades substitution over-editing for word deletion, so it is set
alongside $\wdel$.

\appendix

\section{Generative-function interface weight identities}
\label{app:gfi}

The filter and the rejuvenation move are built on GenJAX's generative-function interface
(\GFI). We recall the identities used above (the genjax/Gen semantics). For a generative function
with density $p(t;a)$ over choice maps $t$:
\begin{description}[leftmargin=1.2em]
  \item[\code{generate}/\code{importance}.] With constraint $c\subseteq t$ and internal proposal
        $q$ over the unconstrained $u$, returns $(t,w)$ with $w = \log p(t;a) - \log q(u\given c;a)$;
        when $c=t$ (fully constrained), $w = \log p(t;a)$. The SMC incremental weight
        \eqref{eq:wstep} is this $w$ for the per-step kernel, whose only stochastic address is the
        proposed word and whose \code{factor} carries $\lse_c S(c)$.
  \item[\code{edit}/\code{Update}.] Overwriting addresses $\mathcal B$ with $u'$ (no-change argdiffs)
        returns weight $w_{\mathrm{upd}} = \log\frac{p(t';a)}{p(t;a)}$ and a backward request that
        restores the old values.
  \item[\code{Rejuvenate} (SMCP3).] With proposal $Q$ and argument map $\phi$, the move draws
        $\zeta\sim Q(\cdot;\phi(t))$, \code{Update}s $\mathcal B\gets\zeta$, and scores the reverse,
        returning $\mathcal W = w_{\mathrm{upd}} + \log Q(t|_{\mathcal B};\phi(t')) -
        \log Q(\zeta;\phi(t))$. Thresholding $\log u < \mathcal W$ gives the MH consumption (A);
        folding $e^{\mathcal W}$ into the weight gives the SMCP3 consumption (B). Both appear in
        \S\ref{sec:rejuv-move}; the inlined \code{\_smcp3\_move} threads per-particle proposal scores
        that the stock \code{Rejuvenate} signature cannot.
\end{description}

\section{Certification by exact enumeration}
\label{sec:certification}

The filter's correctness is established on the toy bigram model, where the candidate set is the entire
vocabulary and the channel marginal is small enough to brute-force
(\code{tests/\allowbreak test\_pairhmm\_exact.py}). The substantive gates that pass are:
(i) \code{caprop\_logZ\_\allowbreak matches\_exact}---the filter's $\log\widehat Z$ equals the exact log-marginal
(minus the deterministic leading-spurious constant); (ii) \code{map\_matches\_exact}---the most
probable (maximum-a-posteriori) intended sentence matches exact enumeration for substitution and
spurious-word cases; and
(iii) \code{posterior\_mass\_\allowbreak matches\_exact}---the full posterior over intended sentences
matches in total variation. A further gate (\code{test\_edit\_types\_\allowbreak recovered\_with\_band})
checks that all four
edit behaviours (substitution; spurious $\to$ shorter; missing $\to$ longer; clean $\to$ unchanged)
are recovered at the production band. Because the toy and Pythia run \emph{identical} inference code
(\S\ref{sec:vmap}), and the multi-token path reduces bit-identically to the certified single-token
path (\S\ref{sec:multitoken}), this certification transfers to the Pythia configuration by
construction; Pythia runs are smoke tests at small $P$, not correctness gates.

\begin{thebibliography}{9}
\bibitem{genjax} MIT Probabilistic Computing Project. \emph{GenJAX}: probabilistic programming with
programmable inference in JAX. \url{https://github.com/genjax-dev/genjax}.
\bibitem{gen} M.~F.~Cusumano-Towner, F.~A.~Saad, A.~K.~Lew, V.~K.~Mansinghka. Gen: a general-purpose
probabilistic programming system with programmable inference. \emph{PLDI}, 2019.
\bibitem{smcp3} A.~K.~Lew, G.~Matheos, et al. SMCP3: Sequential Monte Carlo with probabilistic program
proposals. \emph{AISTATS}, 2023.
\bibitem{chopin2020} N.~Chopin, O.~Papaspiliopoulos. \emph{An Introduction to Sequential Monte Carlo}.
Springer, 2020.
\bibitem{penzai} The Penzai authors. \emph{Penzai}: a JAX research toolkit for neural networks.
\url{https://github.com/google-deepmind/penzai}.
\bibitem{symspell} W.~Garbe. SymSpell: 1 million times faster spelling correction. 2012.
\end{thebibliography}

\end{document}
